% Respuesta a los revisores de BIP2026. Todo número medido va con \result{},
% igual que en main.tex, y verify_paper lo controla:
%   uv run python -c "from ember.paper_sync import verify_paper; print(verify_paper('paper/response.tex') or 'ok')"
\documentclass[11pt]{article}
\usepackage[T1]{fontenc}
\usepackage[margin=2.3cm]{geometry}
\usepackage{amsmath,amssymb}
\usepackage{enumitem}
\usepackage[hidelinks]{hyperref}
% Generado por ember.paper_sync. No editar a mano.
%
% \result{clave}{valor} imprime el valor y deja registrada su procedencia.
% `ember.paper_sync.verify_paper` verifica que el valor coincida con lo medido
% en results/. Un número sin macro es un número que nadie está verificando.
\newcommand{\result}[2]{#2}


\setlist{itemsep=2pt,topsep=3pt}
\setlength{\parindent}{0pt}
\setlength{\parskip}{5pt}
\newcommand{\comment}[1]{\par\medskip\noindent\textbf{#1}\par}

\begin{document}

\begin{center}
{\Large\bfseries Response to the reviewers}\\[4pt]
{\large EMBER: Emergent Memory-Based Encoding and Reactivation}\\[2pt]
BIP2026 --- revised submission
\end{center}

We thank the three reviewers for careful and constructive reports. They
prompted seven new experiments, a rewrite of the bibliography, and a
substantially more precise paper. Several of our claims changed as a result,
and we say so explicitly below rather than only where a reviewer asked.

Section numbers refer to the revised manuscript, which keeps the 8-page limit
including references: I Introduction; II Background and related work; III
Method (A design space, B experimental protocol); IV The e-MDB proxy and the
frontier (A where the proxy lands, B two axes, C how far the advantage holds);
V The regime transition; VI Memory substrates; VII Real data (A CIFAR-100, B
MiniGrid); VIII Discussion and limitations; IX Conclusion. Every measured
number in the paper and in this letter is bound to the result file of the
experiment that produced it and checked automatically before compilation.

\section*{Summary of changes}

\textbf{New experiments.}
\begin{itemize}
\item \textbf{Merge-threshold sensitivity} (R1.4, R2.2): a sweep from 0.30 to 0.97 on synthetic data and CIFAR-100.
\item \textbf{Imperfect salience} (R1.3): overlapping, noisy, delayed and misleading surprise, and novel distractors.
\item \textbf{External baselines and the real e-MDB buffer} (R1.6, R2.3): reservoir sampling, LRU, LFU, a utility cache, coverage selection and prioritized replay, plus the buffer's actual sequential read.
\item \textbf{Robustness of the regime hypothesis} (R1.4, R3.1): 23 stream families, plus 14 held-out families with fresh seeds.
\item \textbf{Extended MiniGrid} (R1.7, R2.1): the rollout as the statistical unit, 200 rollouts, four environments, four policies, five definitions of importance, and a null-signal control.
\item \textbf{Substrate cost} (R1.5): bytes, time, an equal-budget comparison, a counter/trace-list ablation, and SDM sensitivity.
\item \textbf{Downstream behaviour} (R1.9): an episodic controller fed only by the memory.
\end{itemize}

\textbf{What survives.} The central comparison holds:
\begin{itemize}
\item On the synthetic battery, the frontier configuration beats the
nearest-neighbour proxy of the e-MDB buffer:
\result{exp01_nas_full:frontier.score|f3}{0.818} against
\result{exp01_nas_full:incumbent.score|f3}{0.106} aggregate, and
\result{exp01_nas_full:frontier.scores.rare_retention|f3}{0.983} against
\result{exp01_nas_full:incumbent.scores.rare_retention|f3}{0.050} rare-event retention.
\item It keeps the top rank in 13 of 14 held-out stream families.
\item No external replacement rule beats it.
\item The buffer's actual sequential read scores below the proxy.
\item In MiniGrid it retains
\result{exp11_minigrid_extended:exp06_ampliado.aleatoria_recompensa.resultados.frontera.tasa|pct}{47.2}\,\%
of rewarded events against
\result{exp11_minigrid_extended:exp06_ampliado.aleatoria_recompensa.resultados.FIFO.tasa|pct}{0.7}\,\%.
\end{itemize}

\textbf{What changed.} Four results now bound the claims, and each is stated in the abstract or in §VIII:
\begin{enumerate}
\item The regime transition is governed by the \emph{effective} ratio of
consolidated traces to capacity, not by the nominal $K/C$ (R1.4).
\item The conditional salience effect requires a separable surprise signal and
fades as the signal degrades (R1.3).
\item SDM's retention comes from gated eviction with exact erasure, not from
superposition, and at an equal byte budget a FIFO store outperforms it (R1.5).
\item Better retention does not yet translate into better behaviour (R1.9).
\end{enumerate}

\textbf{Corrections we found ourselves.} Integrating the new results, we ran an
independent audit of every empirical claim against the result files. Besides
the points raised by the reviewers, it found the following statements to be
stronger than our data support. All are corrected.
\begin{itemize}
\item \textbf{The two-axes claim (§IV-B).} The submitted version stated that
changing either axis that separates the frontier from the proxy, alone, leaves
the architecture bit-identical to the proxy \emph{in every combination of the
remaining axes}. This is true for the salience gate. It is also true for the
eviction rule in the proxy's own context (append writing, no read
reinforcement), but not in general: once read reinforcement or merge writing
is enabled, weakest-trace eviction alone does change the score. The revised
text restricts the claim to what holds. The conclusion drawn from it is
unaffected: relative to the proxy, neither axis does anything alone.
\item \textbf{The capacity invariance of the crossover.} The three capacities
of the original grid placed the crossover in the same grid cell, which we had
described as ``the same $r$''. The finer sweep shows it moves from
\result{exp10_law_robustness:parte_a.colapso.r_nom.cruces.C005|f2}{0.71} at
$C=5$ to \result{exp10_law_robustness:parte_a.colapso.r_nom.cruces.C160|f2}{0.98}
at $C=160$.
\item \textbf{The MiniGrid retention.} We now report it as
\result{exp11_minigrid_extended:exp06_ampliado.aleatoria_recompensa.resultados.frontera.tasa|pct}{47.2}\,\%
(95\,\% CI
\result{exp11_minigrid_extended:exp06_ampliado.aleatoria_recompensa.resultados.frontera.cluster_low|pct}{39.0}--\result{exp11_minigrid_extended:exp06_ampliado.aleatoria_recompensa.resultados.frontera.cluster_high|pct}{56.3}\,\%)
instead of ``more than half''.
\item \textbf{The CIFAR-100 precondition.} It is now stated once and
consistently, as separability rather than as a reachable threshold (see R2.2).
\end{itemize}

%======================================================================
\section*{Reviewer 1}

\comment{R1.1 --- References [1] and [5] cannot be located; [6] does not support the claim; [8] has an inaccurate title.}
The reviewer is right, and we are grateful this was caught. References [1] and
[5] were our error: they did not correspond to existing publications. They were
not anonymized self-citations; the authors do not belong to the group that
develops e-MDB. We rebuilt the bibliography and verified every entry against a
primary source (DOI via Crossref, PubMed, the publisher's page or the official
proceedings), and checked that each citation supports what it is cited for.
\begin{itemize}
\item \textbf{[1] removed.} e-MDB is now cited through its actual sources from
the Integrated Group for Engineering Research (GII), Universidade da Coruña:
\begin{itemize}
\item Romero's doctoral thesis (UDC, 2022), which presents e-MDB;
\item Romero, Meden, Bellas and Duro, \emph{ICAE} 30(3), 2023, on e-MDB, P-nodes and LOLA;
\item Romero, Bellas and Duro, \emph{Sensors} 23(3), 2023, which defines LOLA;
\item Becerra et al., \emph{Neurocomputing} 452, 2021, on the motivational engine;
\item Duro et al., \emph{IJNS} 29(6), 2019, on the long-term memory;
\item Bellas et al., \emph{IEEE TAMD} 2(4), 2010, on MDB, the predecessor of e-MDB.
\end{itemize}
The description of the episodic buffer as a bounded deque, read by index or by
training batch, appears in no publication. It comes from GII's public
repository (\texttt{episodic\_buffer.py}), which we now cite as software and
identify as such in §II.
\item \textbf{[5] removed.} The prediction-error gate is now supported by
primary literature:
\begin{itemize}
\item Lisman and Grace (2005) and Shohamy and Adcock (2010), on dopaminergic modulation of hippocampal encoding;
\item Rouhani, Norman and Niv (2018), where larger reward prediction errors improve episodic memory;
\item Schultz, Dayan and Montague (1997), for the reward-prediction-error account used in §VII-B.
\end{itemize}
\item \textbf{[6] removed.} It concerned the serial-position effect. The
sentence in §I is rewritten and now rests on Anderson and Schooler (1991),
memory availability tracking need probability, and Richards and Frankland
(2017), forgetting as a regulated process in the service of decision making.
\item \textbf{[8] corrected} to the exact NeurIPS 2023 Datasets and Benchmarks
entry, with the full author list.
\item \textbf{All other entries} now have full authors, exact titles, venues,
volumes, pages and DOIs where they exist. Semon (1904) carries its full title,
and Chaudhry et al.\ (2019) is listed as arXiv/workshop.
\end{itemize}

\comment{R1.2 --- The experimental protocol is not self-contained.}
We added §III-B, ``Experimental protocol'', before any result, with every
detail taken from the code that produced the results:
\begin{itemize}
\item the stream generator: centres, noise, recurrence, and the position and
independence of rare events;
\item the prediction-error distributions of common and rare events;
\item novelty and the gate formulas;
\item the importance tag, which is used only for scoring;
\item the reconstruction gate R1--R4 and its threshold;
\item the battery T1--T3, with query generation and hit criterion;
\item the aggregate score;
\item the number of seeds or rollouts per experiment;
\item the computation of $\eta^2$, interactions and bootstrap intervals (see R1.8).
\end{itemize}

\comment{R1.3 --- Results follow from the construction of the benchmark.}
We agree. The paper now states that, with non-overlapping surprise ranges,
retaining rare events follows almost by construction. We then ran the missing
test: retention under imperfect salience (five seeds). Surprise is degraded by
overlapping ranges, additive noise, delayed credit and partial inversion, and,
on a separate axis, by novel distractors, since novelty also separates rare
events perfectly by construction. The result qualifies our claims, and we
corrected them wherever they appear (abstract, contributions, §IV-B, §IV-C, §VI,
§VIII and §IX).
\begin{itemize}
\item \textbf{The conditional salience ratio depends on the signal.} With the
prediction-error gate alone, it falls from
\result{exp08_imperfect_salience:settings.main.conditions.clean.search.salience_interaction.ratio_vs_constant.pred_error.mean|f2}{16.27}$\times$
on the clean signal to
\result{exp08_imperfect_salience:settings.main.conditions.auc500.search.salience_interaction.ratio_vs_constant.pred_error.mean|f2}{0.80}$\times$
(95\,\% CI
\result{exp08_imperfect_salience:settings.main.conditions.auc500.search.salience_interaction.ratio_vs_constant.pred_error.ci.0|f2}{0.17}--\result{exp08_imperfect_salience:settings.main.conditions.auc500.search.salience_interaction.ratio_vs_constant.pred_error.ci.1|f2}{2.82})
under full overlap. The combined gate keeps more only through novelty. With
25\,\% novel distractors it stays at
\result{exp08_imperfect_salience:settings.main.conditions.distr025.search.salience_interaction.ratio_vs_constant.both.mean|f2}{11.80}$\times$,
and adding chance-level prediction error brings it to
\result{exp08_imperfect_salience:settings.main.conditions.auc500_distr025.search.salience_interaction.ratio_vs_constant.both.mean|f2}{3.10}$\times$.
The paper now attributes the
\result{exp01_nas_full:salience_conditional_ratio|f2}{14.41}$\times$ explicitly
to a perfectly separable signal.
\item \textbf{Against the proxy, the advantage survives every degradation at
$r=0.25$.} Even with surprise at chance, the frontier retains
\result{exp08_imperfect_salience:settings.main.conditions.auc500.policies.frontier.recall.mean|f2}{0.60}
against the proxy's
\result{exp08_imperfect_salience:settings.main.conditions.auc500.policies.fifo.recall.mean|f2}{0.03}.
\item \textbf{Against the best configuration without salience, the advantage
has a signal-quality threshold.} That configuration already retains
\result{exp08_imperfect_salience:settings.main.conditions.clean.policies.no_salience_best.recall.mean|f2}{0.76}
by merging at $r=0.25$. The paired difference excludes zero down to an
empirical AUC of
\result{exp08_imperfect_salience:settings.main.conditions.auc950.auc.mean|f2}{0.95}
with overlap,
\result{exp08_imperfect_salience:settings.main.conditions.noise030.auc.mean|f2}{0.96}
with noise and
\result{exp08_imperfect_salience:settings.main.conditions.mislead010.auc.mean|f2}{0.91}
with 10\,\% inverted signal. It holds to
\result{exp08_imperfect_salience:settings.selection.conditions.auc800.auc.mean|f2}{0.84}
at $r=2$, where merging frees no room. Delay is the most damaging degradation:
shifted by one step, the frontier retains
\result{exp08_imperfect_salience:settings.main.conditions.delay1.policies.frontier.recall.mean|f2}{0.21}.
\item \textbf{SDM's clean-signal retention is qualified.} SDM drops to
\result{exp08_imperfect_salience:settings.main.conditions.auc990.architectures.SDM.recall.mean|f2}{0.81}
at AUC
\result{exp08_imperfect_salience:settings.main.conditions.auc990.auc.mean|f2}{0.99},
and to
\result{exp08_imperfect_salience:settings.main.conditions.auc950.architectures.SDM.recall.mean|f2}{0.60}
at AUC 0.95, below ENN
(\result{exp08_imperfect_salience:settings.main.conditions.auc950.architectures.ENN.recall.mean|f2}{0.75})
in that stream.
\item \textbf{Precision and recall.} With memory full and as many rare events
as slots, precision equals recall arithmetically. In the $r=2$ stream,
precision is bounded at 0.5 and tracks recall. The paper says so in §III-B
rather than reporting a redundant metric.
\end{itemize}

\comment{R1.4 --- Calling it a ``threshold law'' is premature.}
We agree, and the paper now states it as a \emph{regime hypothesis} supported
under the evaluated conditions. §V notes that the existence of a transition is
close to a restatement of what capacity means; what the grid adds is that the
switch is abrupt. We then ran the robustness test the reviewer asked for, and
it changed the statement.
\begin{itemize}
\item \textbf{Nominal $r$ is not the governing variable.} We repeated the
sweep on 23 stream families, each moving one knob: Zipf frequencies, drift,
intra-prototype noise, centre similarity, rare-event prevalence, recurrence,
and $C$ from 5 to 160. The nominal crossover ranges from
\result{exp10_law_robustness:parte_a.colapso.r_nom.cruce_min|f2}{0.16} to
\result{exp10_law_robustness:parte_a.colapso.r_nom.cruce_max|f2}{1.32}
(s.d.\ of its log
\result{exp10_law_robustness:parte_a.colapso.r_nom.sd_log|f2}{0.46}): drift
and noise pull it far below 1.
\item \textbf{The effective ratio is.} Over $r_{eff}=K_{effective}/C$, the
traces left after consolidation, crossovers collapse to
[\result{exp10_law_robustness:parte_a.colapso.r_eff.cruce_min|f2}{0.56},
\result{exp10_law_robustness:parte_a.colapso.r_eff.cruce_max|f2}{1.28}]
(s.d.\ \result{exp10_law_robustness:parte_a.colapso.r_eff.sd_log|f2}{0.17}).
That is tighter than nominal $r$ in
\result{exp10_law_robustness:parte_a.colapso.r_eff.frac_bootstrap_mas_colapsada|pct}{99.8}\,\%
of bootstrap replicates, and tighter than four other pressure measures fixed
in advance. In the original generator $r_{eff}=r$, so the main figure is
unchanged. The hypothesis is now stated over $r_{eff}$ throughout.
\item \textbf{The merge threshold does not place the crossover.} Sweeping it
from 0.30 to 0.97 leaves the crossover at
$r\approx\result{exp07_merge_sensitivity:umbral_fusion.synthetic.umbrales.t085.cruce_r_nominal|f2}{0.87}$
for every threshold from 0.40 to 0.85. Over $K_{effective}$ it stays between
\result{exp07_merge_sensitivity:umbral_fusion.synthetic.umbrales.t040.cruce_r_efectivo|f2}{0.86}
and
\result{exp07_merge_sensitivity:umbral_fusion.synthetic.umbrales.t030.cruce_r_efectivo|f2}{1.02}
from 0.30 to 0.90. A sweep of the gate's gain leaves the pattern unchanged.
\end{itemize}

\comment{R1.5 --- Substrates compared under non-equivalent storage and compute budgets.}
We measured this, with ten seeds and bootstrap intervals, and the result
qualifies the paper (§VI, Table III).
\begin{itemize}
\item \textbf{Cost.} At $C=20$, SDM occupies
\result{exp12_substrate_cost:contabilidad.razon_vs_FIFO_C20.SDM|f2}{41.96}$\times$
the bytes of FIFO and ENN. A read takes
\result{exp12_substrate_cost:tiempos.C20_d32.SDM.read_us|int}{28}\,$\mu$s
against \result{exp12_substrate_cost:tiempos.C20_d32.FIFO.read_us|int}{12}\,$\mu$s.
Table III now has a KiB column.
\item \textbf{Equal budget.} With SDM's footprint, a FIFO store holds
\result{exp12_substrate_cost:igual_presupuesto.defaults.B147456.filas.FIFO.capacity|int}{921}
traces and outperforms it on the battery,
\result{exp12_substrate_cost:igual_presupuesto.defaults.B147456.filas.FIFO.battery_mean|f3}{0.924}
against
\result{exp12_substrate_cost:igual_presupuesto.defaults.B147456.filas.SDM.battery_mean|f3}{0.745}.
At that budget the stream no longer pressures capacity, and the paper says so.
\item \textbf{Counters versus trace list.} The trace list alone, read by
nearest neighbour, keeps every rare event and scores
\result{exp12_substrate_cost:ablacion_sdm.variantes.SDM-b.battery_mean|f3}{0.788}
against \result{exp12_substrate_cost:ablacion_sdm.variantes.SDM.battery_mean|f3}{0.741}
for the full SDM. Counters alone keep
\result{exp12_substrate_cost:ablacion_sdm.variantes.SDM-a.battery.rare_retention.mean|f3}{0.005}
of rare events if they cannot subtract an evicted trace, but
\result{exp12_substrate_cost:ablacion_sdm.variantes.SDM-a+.battery.rare_retention.mean|f3}{1.000}
with exact erasure. Retention therefore needs gated eviction with exact
erasure, not superposition. We corrected the sentence that attributed
retrieval to superposition, and the claim that SDM was the only substrate
combining reconstruction and retention.
\item \textbf{Sensitivity.} The battery mean moves
\result{exp12_substrate_cost:sensibilidad_sdm.rango_battery_n_hard|f3}{0.091}
between 64 and 2048 hard locations. Rare-event retention is perfect for
activation fractions from 0.005 to 0.2 and drops to
\result{exp12_substrate_cost:sensibilidad_sdm.activation_frac.f0p5.battery.rare_retention.mean|f3}{0.395}
only at 0.5, a setting that also fails the reconstruction gate.
\end{itemize}
Write times, asymptotic complexity and the scaling of hard locations to each
budget are in the released results; they are not in the paper for space.

\comment{R1.6 --- The ``e-MDB FIFO'' is an optimistic proxy.}
We renamed the baseline the \emph{FIFO-NN proxy of the e-MDB buffer}
throughout. The paper states that the deployed buffer retrieves by position or
batch, not by content, and that we chose nearest-neighbour retrieval to isolate
eviction and weighting.

We also evaluated the real buffer (§IV-C). We re-implemented its read as in
GII's public code: a scan in insertion order that returns the first episode
with cosine $\geq 0.85$. This translation of a query is an assumption, since
the buffer answers no queries. With the same eviction it keeps exactly the
same traces as the proxy, but it fails the reconstruction gate,
\result{exp09_external_baselines:exp01_cell.emdb_sequential.gate_mean.mean|f3}{0.368}
against \result{exp09_external_baselines:exp01_cell.fifo_nn_proxy.gate_mean.mean|f3}{0.903}.
It scores
\result{exp09_external_baselines:exp01_cell.emdb_sequential.battery_mean.mean|f3}{0.010}
and ranks last of
\result{exp09_external_baselines:rank_in_space.emdb_sequential.rank_optimistic|int}{577}.
The gate failure depends on the assumed match criterion: at 0.5 the gate
recovers
(\result{exp09_external_baselines:emdb_threshold_sensitivity.tau_050.gate_mean.mean|f3}{0.826}).
The score does not
(\result{exp09_external_baselines:emdb_threshold_sensitivity.tau_050.battery_mean.mean|f3}{0.067}),
and the paper states only the latter as robust. The proxy's rank is therefore
an upper bound for the deployed buffer.

\comment{R1.7 --- MiniGrid is limited; importance is defined by the task; Wilson intervals treat events as independent.}
We addressed this with a new experiment (§VII-B, and §VIII for delayed
importance).
\begin{itemize}
\item \textbf{Statistical unit.} With the rollout as the unit (cluster
bootstrap) on the original 40 rollouts, the frontier retains
\result{exp06_minigrid:comparacion_aleatoria_recompensa.resultados.frontera.tasa|pct}{52.6}\,\%,
CI
[\result{exp11_minigrid_extended:exp06_original.aleatoria_recompensa.resultados.frontera.cluster_low|pct}{37.9},
\result{exp11_minigrid_extended:exp06_original.aleatoria_recompensa.resultados.frontera.cluster_high|pct}{69.0}]\,\%,
and FIFO
\result{exp06_minigrid:comparacion_aleatoria_recompensa.resultados.FIFO.tasa|pct}{2.6}\,\%,
CI
[\result{exp11_minigrid_extended:exp06_original.aleatoria_recompensa.resultados.FIFO.cluster_low|pct}{0.0},
\result{exp11_minigrid_extended:exp06_original.aleatoria_recompensa.resultados.FIFO.cluster_high|pct}{8.6}]\,\%.
The design effect is
\result{exp11_minigrid_extended:exp06_original.aleatoria_recompensa.resultados.frontera.deff|f2}{0.92},
so the conclusion does not change. The paper now reports rollout-level
intervals.
\item \textbf{More rollouts.} With 200 rollouts in five independent seed
blocks, the frontier retains
\result{exp11_minigrid_extended:exp06_ampliado.aleatoria_recompensa.resultados.frontera.tasa|pct}{47.2}\,\%
[\result{exp11_minigrid_extended:exp06_ampliado.aleatoria_recompensa.resultados.frontera.cluster_low|pct}{39.0},
\result{exp11_minigrid_extended:exp06_ampliado.aleatoria_recompensa.resultados.frontera.cluster_high|pct}{56.3}]\,\%
against
\result{exp11_minigrid_extended:exp06_ampliado.aleatoria_recompensa.resultados.FIFO.tasa|pct}{0.7}\,\%
for FIFO.
\item \textbf{Null-signal control.} The same frontier fed a uniform random
signal retains
\result{exp11_minigrid_extended:matriz.MemoryS13.aleatoria.recompensa.senales.azar.frontera.tasa|pct}{0.7}\,\%.
The gain comes from the signal, not from strength-based eviction.
\item \textbf{Importance not defined by the task.} We crossed four
environments (MemoryS13, DoorKey-6x6, FourRooms, KeyCorridorS3R2), four
policies (uniform, forward-biased, noisy BFS planner, tabular Q-learning), five
definitions of importance and five candidate signals, for
\result{exp11_minigrid_extended:resumen.n_celdas_total|int}{350} cells. The
gain over the null control tracks the signal's AUC (Spearman
\result{exp11_minigrid_extended:resumen.spearman_auc_ganancia|f2}{0.67}):
\result{exp11_minigrid_extended:resumen.por_auc.t5.ganancia_vs_azar_media|pct}{37.3}
points where AUC $\geq 0.97$, and
\result{exp11_minigrid_extended:resumen.por_auc.t1.ganancia_vs_azar_media|pct}{1.8}
near chance. Count-based novelty beats the null control in
\result{exp11_minigrid_extended:resumen.novedad_estado.count_novelty.gana_vs_azar|int}{16}
of \result{exp11_minigrid_extended:resumen.novedad_estado.count_novelty.n_celdas|int}{16}
cells when the important states are the novel ones.
\item \textbf{Delayed importance remains open.} No causal signal retains the
steps preceding a reward. Reward prediction error beats the null control in
\result{exp11_minigrid_extended:resumen.previa_k.reward_pe.gana_vs_azar|int}{0}
of \result{exp11_minigrid_extended:resumen.previa_k.reward_pe.n_celdas|int}{15}
cells and TD error in
\result{exp11_minigrid_extended:resumen.previa_k.td_error.gana_vs_azar|int}{2}.
Only the discounted return, which is not causal, succeeds consistently
(\result{exp11_minigrid_extended:resumen.previa_k.retorno.gana_vs_azar|int}{15}
of 15). §VIII now states this as the open problem, with evidence.
\end{itemize}

\comment{R1.8 --- Statistics are insufficiently explained.}
§III-B now defines the statistics we use.
\begin{itemize}
\item \textbf{$\eta^2$} is SS$_{between}$/SS$_{total}$ over the 576
configurations, a one-way marginal decomposition. In a full balanced factorial
the main effects are orthogonal.
\item \textbf{Interactions} are two-way cell-mean residuals, reported
separately and never added to an axis's $\eta^2$.
\item \textbf{There is no error term or significance test,} because results
are deterministic given the seed.
\item \textbf{Bootstrap intervals} use 2000 percentile replicates with the
seed as the unit.
\end{itemize}
A paragraph on structural dependence explains that strength and reinforcement
are inert under three of the four eviction rules. Main effects are therefore
read as attribution over the whole space, and conditional effects are reported
where the dependence is known. The interaction shares quantify it:
strength$\times$evict
\result{exp01_nas_full:variance_share.strength×evict|pct}{4.3}\,\% and
write$\times$evict
\result{exp01_nas_full:variance_share.write×evict|pct}{7.8}\,\%, against
\result{exp01_nas_full:main_effects.3.eta2|pct}{1.4}\,\% for the main effect of
strength.

\comment{R1.9 --- Retention is not an effect on learning; no integration into e-MDB; non-biological literature is missing.}
\begin{itemize}
\item \textbf{Effect on behaviour} (§VIII, abstract and conclusion). An
episodic controller whose only experience is the memory faces four MiniGrid
goal variants that alternate and recur (ten seeds, paired bootstrap).
\begin{itemize}
\item \textbf{At $C=300$:} the frontier earns more than FIFO on the first
episodes of a returning task
(+\result{exp13_downstream_learning:resumen.C300.diferencias_pareadas.frontera_menos_FIFO.reaparicion.media|f3}{0.038},
CI
[\result{exp13_downstream_learning:resumen.C300.diferencias_pareadas.frontera_menos_FIFO.reaparicion.ci_low|f3}{0.005},
\result{exp13_downstream_learning:resumen.C300.diferencias_pareadas.frontera_menos_FIFO.reaparicion.ci_high|f3}{0.071}]),
but it integrates less return over the stream: area under the curve
\result{exp13_downstream_learning:resumen.C300.diferencias_pareadas.frontera_menos_FIFO.auc.media|f3}{-0.075},
CI
[\result{exp13_downstream_learning:resumen.C300.diferencias_pareadas.frontera_menos_FIFO.auc.ci_low|f3}{-0.098},
\result{exp13_downstream_learning:resumen.C300.diferencias_pareadas.frontera_menos_FIFO.auc.ci_high|f3}{-0.053}].
\item \textbf{At $C=1000$:} the pattern repeats
(+\result{exp13_downstream_learning:resumen.C1000.diferencias_pareadas.frontera_menos_FIFO.reaparicion.media|f3}{0.109}
and
\result{exp13_downstream_learning:resumen.C1000.diferencias_pareadas.frontera_menos_FIFO.auc.media|f3}{-0.064}).
\item \textbf{With decay:} adding it, from the same space, brings the area to
the FIFO level, but the returning-task advantage then shrinks to within noise
at $C=300$ and is lost at $C=1000$.
\end{itemize}
No tested configuration is significantly better than FIFO on both measures.
The frontier does beat reservoir sampling on integrated return. We state this
plainly: better retention is not yet better behaviour. Its effect on the
models e-MDB trains from its buffer remains untested.
\item \textbf{Integration.} §VIII states that EMBER is not integrated into
the running e-MDB, and that Fig.~1 is a design, not an implementation.
\item \textbf{Non-biological memory management} (§II):
\begin{itemize}
\item LRU-K and ARC;
\item reservoir sampling;
\item Chaudhry et al.;
\item iCaRL and GSS;
\item selective experience replay;
\item experience replay for continual learning.
\end{itemize}
These methods are now also compared empirically (see R2.3).
\end{itemize}

\comment{R1.10 --- Release code and configurations during review.}
The ``Code and data availability'' section states that the code,
configurations, results and provenance manifests will be released in a public
repository under an open licence. Every number in the paper is generated from
those results by the released code.

\comment{R1.11 --- Recommended improvements.}
Each recommendation is addressed above:
\begin{itemize}
\item consolidated protocol (R1.2);
\item hypothesis framing and more generators (R1.4);
\item imperfect signals, with precision and recall (R1.3);
\item held-out families and external baselines (R1.4, R2.3);
\item equivalent substrate budgets and hyperparameter sensitivity (R1.5, R2.2);
\item rollout-level MiniGrid statistics, more environments, and the effect on learning (R1.7, R1.9);
\item comparison against the real sequential buffer (R1.6).
\end{itemize}
Full integration into a running e-MDB is declared as future work.

%======================================================================
\section*{Reviewer 2}

\comment{R2.1 --- Scope of validation: one dataset, one extractor, one environment, no physical robot.}
MiniGrid now spans four environments and four policies, two of them directed
(a noisy BFS planner and tabular Q-learning). It uses 200 rollouts with the
rollout as the unit (§VII-B; details in R1.7). Salience transfers when the
signal is aligned with what matters and fails for delayed importance. §VIII
states the limits no experiment here resolves: discrete grid worlds, a single
visual extractor, no physical robot, and no integration into the running
e-MDB.

\comment{R2.2 --- Dependence on the 0.85 merge threshold.}
The reviewer was right that this needed testing. We swept the threshold from
0.30 to 0.97 in both domains. The CIFAR-100 result does not depend on the 0.85
calibration.
\begin{itemize}
\item \textbf{CIFAR-100} (§VII-A). Intra-class similarity (median
\result{exp07_merge_sensitivity:umbral_fusion.cifar100.similitudes.intra.p50|f2}{0.38})
overlaps inter-class similarity (95th percentile
\result{exp07_merge_sensitivity:umbral_fusion.cifar100.similitudes.inter.p95|f2}{0.42}),
so no threshold yields merges that are both frequent and pure. At 0.30 the
write axis regains dominance at low $r$, but not as useful compression:
\begin{itemize}
\item only \result{exp07_merge_sensitivity:umbral_fusion.cifar100.umbrales.t030.pureza_fusion_r_min|pct}{47.7}\,\%
of merges join the same class;
\item \result{exp07_merge_sensitivity:umbral_fusion.cifar100.umbrales.t030.raros_absorbidos_r_min|pct}{70.0}\,\%
of rare events are absorbed on arrival;
\item the best merging configuration retains
\result{exp07_merge_sensitivity:umbral_fusion.cifar100.umbrales.t030.celdas.r0p25.puntaje_mejor_merge.media|f2}{0.30}
against
\result{exp07_merge_sensitivity:umbral_fusion.cifar100.umbrales.t030.celdas.r0p25.puntaje_frontera.media|f2}{1.00}
for the append frontier.
\end{itemize}
The viability criterion we fixed before running is met by no threshold: merge
rate $\geq 0.25$, purity $\geq 0.9$, and a dominant write axis at low $r$.
\item \textbf{Correction.} The precondition is no longer ``intra-prototype
similarity above the merge threshold'' but \emph{separability}: a threshold
that repeated encounters clear and distinct ones do not. Separability was
measured with class labels. An online, label-free estimator is untested, and
the paper says so. We withdrew the sentence suggesting that a per-domain
threshold could change the outcome; a fine-tuned extractor might.
\item \textbf{Synthetic} (§V). The crossover is robust to the threshold (see
R1.4).
\end{itemize}

\comment{R2.3 --- More memory-management baselines.}
§II now covers the non-biological literature: caching, reservoir sampling and
replay selection. §IV-C compares against it empirically. Seven external rules
each replace only the proxy's eviction, on T1--T3 with five seeds:
\begin{itemize}
\item the frontier scores
\result{exp09_external_baselines:exp01_cell.frontier.battery_mean.mean|f3}{0.824};
\item LRU, LFU and a utility cache stay at the proxy's level
(\result{exp09_external_baselines:exp01_cell.lru.battery_mean.mean|f3}{0.106},
\result{exp09_external_baselines:exp01_cell.lfu.battery_mean.mean|f3}{0.096},
\result{exp09_external_baselines:exp01_cell.utility_cache.battery_mean.mean|f3}{0.106});
\item reservoir sampling reaches
\result{exp09_external_baselines:exp01_cell.reservoir.battery_mean.mean|f3}{0.306};
\item coverage selection reaches
\result{exp09_external_baselines:exp01_cell.coverage.battery_mean.mean|f3}{0.439};
\item stochastic prioritized replay reaches
\result{exp09_external_baselines:exp01_cell.per_stoch.battery_mean.mean|f3}{0.500};
\item greedy surprise-prioritized eviction reaches
\result{exp09_external_baselines:exp01_cell.per_min.battery_mean.mean|f3}{0.789}
(rank \result{exp09_external_baselines:rank_in_space.per_min.rank_optimistic|int}{4}--\result{exp09_external_baselines:rank_in_space.per_min.rank_pessimistic|int}{5}).
\end{itemize}
Greedy surprise-prioritized eviction coincides exactly with a configuration
already in the space. The frontier beats it by
\result{exp09_external_baselines:advantage_of_frontier.per_min.battery_mean.mean|f3}{0.035}
[\result{exp09_external_baselines:advantage_of_frontier.per_min.battery_mean.ci.0|f3}{0.028},
\result{exp09_external_baselines:advantage_of_frontier.per_min.battery_mean.ci.1|f3}{0.043}].
The paper says this plainly: the frontier's advantage comes from reading
surprise at eviction, a principle prioritized replay already uses, not from a
mechanism unique to bioinspired memory.

%======================================================================
\section*{Reviewer 3}

\comment{R3.1 --- Scope of validation and generalization to robots.}
See R2.1 for the environments and the stated limits, and R1.4 for the regime
hypothesis, which is no longer stated as a law. On generalization beyond the
generator: the frontier selected on the standard stream, re-scored on 14
held-out families with fresh seeds, keeps rank interval [1,\,1] in all but
one, and
[\result{exp10_law_robustness:parte_b.familias.raros_01pct.bateria.frontera_rango_optimista|int}{15},\,\result{exp10_law_robustness:parte_b.frontera_peor_rango_pesimista_bateria|int}{20}]
with 0.1\,\% rare events. The ranking of all 576 configurations correlates
with the standard one at Spearman
$\geq$\,\result{exp10_law_robustness:parte_b.spearman_min_bateria|f2}{0.77}.

\comment{R3.2 --- A diagram of the EMBER/e-MDB integration.}
We added Fig.~1 (§II). It shows:
\begin{itemize}
\item e-MDB's long-term memory and its cognitive nodes, and the motivational engine;
\item the path of each episode into EMBER, in place of the FIFO episodic buffer;
\item the write path, with the gate fed by the surprise of e-MDB's own models and weakest-trace eviction;
\item the content-based read path that serves model training and deliberation.
\end{itemize}
The caption states that it is a design, not an implemented integration.

\end{document}
